\documentclass[11pt]{article}
\usepackage[T1]{fontenc}
\usepackage[utf8]{inputenc}
\usepackage{lmodern,amsmath,amssymb,amsthm}
\usepackage[margin=1in]{geometry}
\usepackage[hidelinks]{hyperref}
\newtheorem{theorem}{Theorem}
\newtheorem{lemma}[theorem]{Lemma}
\newtheorem{corollary}[theorem]{Corollary}
\newcommand{\Z}{\mathbb Z}
\newcommand{\N}{\mathbb N}
\newcommand{\R}{\mathbb R}
\newcommand{\poly}{\operatorname{poly}}
\newcommand{\Nil}{\operatorname{Nil}}
\title{Tautification and the shifted nilcharacter construction}
\author{Project proof draft: /root/astra\_v14\_baseline\_theta}
\date{7 November 2026}
\begin{document}
\maketitle
\noindent\textbf{Status.} This is an author proof proposal. Its internal
geometric input is the rational-subgroup and quotient argument in the separately
frozen local Lemma 13.2 proof \cite{LP13}, which is also awaiting independent
acceptance. The present unit does not assert that the complete inverse theorem
has been repaired.

\section{The exact change of filtration}

All Lie groups and rationality conventions are those of
\cite[Sections 3--6]{GTZ}. A rational subgroup is connected, and its intersection
with the ambient lattice is cocompact. The three domain filtrations used here
are the degree filtration on $\Z^a$, the multidegree filtration on $\Z^a$,
and the degree-rank filtration on $\Z$ from \cite[Definition 6.17]{GTZ}.
For limit objects, defining relations are interpreted by their natural
ultrapower extension as in \cite[Appendix A.2--A.3, pp.~1324--1325]{GTZ}.
Thus vertical equivariance quantifies over every limit point of the
nilmanifold and every limit element of the top group, using the extension
of the standard vertical character. We do not infer this exact internal
identity from an identity asserted only at standard points.
For these three cases respectively put
\[
 H=G_1,\qquad H=\langle G_{e_1},\ldots,G_{e_a}\rangle_{\R},
 \qquad H=G_{(1,0)}.
\]
In the middle case the connected real generated subgroup equals the group
generated by the connected subgroups: the latter is the connected analytic
subgroup with the generated Lie algebra, and such subgroups in a simply
connected nilpotent group are closed. The rational-subgroup lemma in
\cite[Section 2]{LP13} shows that $H$ is rational.

\begin{theorem}\label{thm:taut}
Suppose $G/\Gamma$ is a rational filtered nilmanifold with finite downset of
degrees $J$ for one of these three domain filtrations. Every polynomial
nilsequence $F(g(n)\Gamma)$ admits an exactly equal representation on
$H/(H\cap\Gamma)$ with filtration
\[
 G'_0=H,\qquad G'_i=G_i\quad(i>0).
\]
For degree-rank, $0$ denotes $(0,0)$. This filtration is taut and has the
same positive-degree groups and no larger degree support, hence the same
allowed degree bound. Its test-function Lipschitz
bound depends only on the original fixed filtered nilmanifold and original
Lipschitz bound, not on the polynomial coefficients. For a nilcharacter
of positive degree $d$, the same standard vertical character on $G_d$ is
preserved. These assertions hold uniformly for limit maps and families of
polynomial coefficients with fixed geometric data.
\end{theorem}
\begin{proof}
As in \cite[Section 3]{LP13}, first restrict to $G_0$ and its lattice if
necessary. Choose $g(0)=q\gamma$ with $q$ in a fixed compact representative
set, and set
\[
 h(n)=q^{-1}g(n)\gamma^{-1},\qquad F_q(x\Gamma)=F(qx\Gamma).
\]
Then $h(0)=1$, $F_q(h(n)\Gamma)=F(g(n)\Gamma)$, and $h$ is polynomial
for the old filtration. Indeed every $G_i$ is normal in $G_0$; use the
constant map $g(0)^{-1}$, product closure in \cite[Corollary B.4]{GTZ},
and filtered conjugation by $q$. The Lipschitz bound is uniform because
$q$ lies in a fixed compact set.

We show that $h(n)$ belongs to $H$. In the degree case the difference
$h(n)h(0)^{-1}$ is in $G_1$ for every $n\in\Z^a$, by the derivative
condition. In the degree-rank case the same statement holds with
$G_{(1,0)}$. In the multidegree case connect $0$ to $n$ by a finite path
of positive or negative coordinate steps. At a step in the $j$th coordinate,
$h(v\pm e_j)h(v)^{-1}\in G_{e_j}$. Multiplying these increments along the
path and using $h(0)=1$ proves the assertion. This argument applies to every
integer input and puts no bound on its size.

Every original positive group is contained in $H$. For degree and
degree-rank this follows from nesting and the least positive index.
For multidegree, any positive index dominates some $e_j$.
Consequently the displayed new groups are well-defined subgroups of $H$,
are descending, and satisfy the commutator axioms inherited from $G$,
including those with index zero because $H\subseteq G_0$.
Their intersections with $H\cap\Gamma$ equal their intersections with
$\Gamma$ and are lattices. In degree-rank all the required equalities
$G_{(d,0)}=G_{(d,1)}$ for $d>0$ remain unchanged. The new group at zero is
exactly the required positive group or group generated by the minimal
positive groups, so the new filtration is taut.

All positive iterated derivatives of $h$ have the same required group
membership as before, and the zeroth-order membership was just checked.
Thus $h$ is polynomial for the new filtration. Restrict $F_q$ through the
fixed smooth inclusion $H/(H\cap\Gamma)\longrightarrow G/\Gamma$.
The inclusion is an embedding of compact manifolds, so this costs only a
fixed Lipschitz comparison factor and gives the asserted exact sequence.

For a nilcharacter let $\eta:G_d\longrightarrow\R$ be its standard
vertical character. The proof in \cite[Section 3]{LP13} applies verbatim
with $G_d$ in place of its top cube group: for each $t\in G_d\cap\Gamma$,
the continuous map $q\mapsto\eta(qtq^{-1})$ on connected $G_0$ is
integer-valued by the unit-norm vertical identity at $q\Gamma$.
It is therefore constant. The lattice spans $G_d$ as a connected real
group, so $\eta(qtq^{-1})=\eta(t)$ for all $t\in G_d$ and $q\in G_0$.
Hence $F_q$ and its restriction have the same standard vertical character.

For families, choose one bounded standard representative section of the
fixed compact quotient. Such a section can be formed from finitely many
local smooth sections, taking the first available chart. Apply this one
section to each constant coefficient. The resulting choices are limit maps
when the coefficients are limit maps. Discontinuities in the parameter do
not enter a Lipschitz bound in the orbit variable. All metric constants and
groups remain fixed. The exact identities can be imposed on representatives
on an ultrafilter-large set by the witness argument of
\cite[Section 7]{LP13}; the ordinary proof then passes to the limit.
\end{proof}

This also proves equality of the classes of nilsequences under the old and
taut conventions, at each finite degree support: one inclusion is the theorem
and the reverse inclusion is immediate. It therefore preserves the
equivalence relation defined by a lower-degree nilsequence. It is not merely
an assertion that the top-degree function has a similar replacement.

\begin{corollary}\label{cor:central}
In a taut degree or degree-rank filtration of degree at most $(d,r)$ or $d$,
the ambient group has nilpotency class at most $d$. In a taut multidegree
filtration supported on a finite downset $J$, its class is at most
$D=\max_{a\in J}|a|$. If the support is contained in the downset below one
positive index $d$, the group $G_d$ is central.
\end{corollary}
\begin{proof}
For degree and degree-rank, brackets of $j$ elements of the ambient first
positive group belong respectively to $G_j$ or $G_{(j,0)}$, which vanish
when $j>d$. For multidegree, expand an iterated Lie bracket using the
generating Lie algebras $\mathfrak g_{e_i}$. A bracket of length $j$ lies in
some $\mathfrak g_a$ with $|a|=j$, hence vanishes for $j>D$.
The generated Lie algebra has the asserted nilpotency class, and so does
its simply connected group. Finally $[G_d,G_{e_i}]\subseteq G_{d+e_i}=1$
in the multidegree case. The same argument uses the first positive group
in the other two cases. Since these groups generate the ambient group,
$G_d$ is central.
\end{proof}

\section{A complete shifted nilcharacter construction}

\begin{theorem}\label{thm:shift}
For any of the three domain filtrations above, let $\chi$ be a nilcharacter
of positive degree $d$. Let $i>0$ and let $u$ belong to the domain group at
index $i$. Then
\[
 n\longmapsto\chi(n+u)\otimes\overline{\chi(n)}
\]
is a nilsequence with support strictly below $d$. The representing geometric
data and Lipschitz bound can be fixed independently of $u$ within this
domain group. For limit $u$ the representation and test function can be
chosen in a limit fashion. This supplies the shift assertion required in
\cite[Lemma E.8(iv)]{GTZ} for the three domains used in the proof.
\end{theorem}
\begin{proof}
By Theorem~\ref{thm:taut} represent the same $\chi$ on a taut nilmanifold
$G/\Gamma$, with the same degree and vertical character. Write its
polynomial map as $g$. In the degree-rank case use $i=(1,0)$: any
nonzero element of a positive domain group requires this index, and if the
given domain group is trivial then $u=0$ also belongs to this first group.
Put $U(v)=(v,1)$ and $D(v)=(v,v)$ in $G^2$.
For each index $j$ let $L_j$ be the connected subgroup generated by
\[
 U(G_{j+i})\quad\hbox{and}\quad D(G_j).
\]
These groups are rational by \cite[Lemma 2]{LP13}, and are descending.
The generator commutators have the following memberships:
\begin{align*}
 [D(G_j),D(G_k)]&\subseteq D(G_{j+k}),\\
 [D(G_j),U(G_{k+i})]&\subseteq U(G_{j+k+i}),\\
 [U(G_{j+i}),U(G_{k+i})]&\subseteq U(G_{j+k+2i})
                                      \subseteq U(G_{j+k+i}).
\end{align*}
Passing to generated Lie algebras and using BCH, as in
\cite[Section 4]{LP13}, proves $[L_j,L_k]\subseteq L_{j+k}$, including
index zero. Thus $L_0/\Delta$, with $\Delta=L_0\cap\Gamma^2$, is a rational
filtered nilmanifold of degree at most $d$. In the degree-rank case its
additional equalities hold: $L_{(a,0)}=L_{(a,1)}$ follows from
$G_{(a,0)}=G_{(a,1)}$ and $G_{(a+1,0)}=G_{(a+1,1)}$ for $a>0$.

The difference $a_u(n)=g(n+u)g(n)^{-1}$ is polynomial for the shifted
filtration $(G_{j+i})_j$. Its zeroth-order value belongs to $G_i$.
Every further iterated derivative is an iterated derivative of $g$ with
the derivative in direction $u$ innermost, so the defining index sum has
the additional term $i$. Therefore $U(a_u(n))$ and $D(g(n))$ are polynomial
for $(L_j)$. By \cite[Corollary B.4]{GTZ}, so is their product
\[
 P_u(n)=(g(n+u),g(n)).
\]
This verifies the actual orbit, not only its coordinate projections.

As $i>0$, $G_{d+i}=1$ and $L_d=D(G_d)$. Corollary~\ref{cor:central}
shows that $G_d$ is central in $G_0$, hence $L_d$ is central in $L_0$.
The function
\[
 T(x\Gamma,y\Gamma)=F(x\Gamma)\otimes\overline{F(y\Gamma)}
\]
restricted to $L_0/\Delta$ is invariant under $D(G_d)$, since its two
vertical characters cancel. The rational central-quotient argument in
\cite[Sections 2 and 6]{LP13} therefore gives a nilsequence on
$L_0/L_d$ with the exact value $T(P_u(n)\Delta)$. Its filtration group at
$d$ is trivial and all groups outside the downset below $d$ were already
trivial. Its support is consequently strictly below $d$.
The same argument gives a fixed Lipschitz bound on the quotient.

The quotient may not yet be taut. Apply Theorem~\ref{thm:taut} to this
lower-degree nilsequence, preserving its exact values without enlarging its
degree support.
The groups $L_j$ and their quotient depend on the index $i$ but not on the
element $u$. The map $u\mapsto P_u(0)$ is a limit map when $g$ and $u$ are
limit data. The one fixed bounded section used in Theorem~\ref{thm:taut}
then gives a limit choice of the final normalized representation, with
uniform Lipschitz bounds independent of $u$. This proves all the assertions.
\end{proof}

For the erratum's needed case the domain is $\Z^2$ with multidegree,
$d=(1,s-1)$, $i=(1,0)$ and $u=(k,0)$. The resulting support is contained in
\[
 \{a:a\leq(0,s-1)\}\ \cup\ \{a:a\leq(1,s-2)\},
\]
uniformly in $k$. This is exactly the support needed before applying the
published splitting and vertical Fourier lemmas in the proof of the new
difference decomposition. Those applications, and the subsequent two
Cauchy--Schwarz steps, are separate proof obligations.

\section{Scope of the correction}

The tautification theorem supplies the general representation repair
suggested, but not fully proved, in Section 2 of the author erratum.
It also proves the top centrality and nilpotency-class assertions for the
taut conventions. The shifted construction above supplies the needed
instance of E.8(iv) without leaving its auxiliary rationality or central
quotient implicit. This unit does not claim to have audited every filtration
introduced in the published proof. In particular, E.8(v), E.9, Lemma 12.1,
the Section 12 quotient, and the final modified Proposition 7.3 argument
still need their own exact applications or replacements.

\begin{thebibliography}{9}
\bibitem{GTZ} B. Green, T. Tao and T. Ziegler,
\emph{An inverse theorem for the Gowers $U^{s+1}[N]$-norm},
Annals of Mathematics \textbf{176} (2012), 1231--1372.
\bibitem{LP13} Project proof, \emph{A local replacement for the filtration
argument in GTZ Lemma 13.2}, draft 002, 7 November 2026.
Manifest SHA-256:
\par\noindent{\small\texttt{95771803bed9d9135430152e2e20ad450c46aaae3ef0eeba0cab33bee59fd06c}}
\par\noindent
This is an explicitly linked author proof awaiting independent acceptance.
\end{thebibliography}
\end{document}
